%% ============================================================
%% sections/a5-evidence.tex
%% H. R. 9510 Bill v3.0 - APPENDIX A. VISUAL ENGINEERING EVIDENCE (NON-OPERATIVE).
%% The visualization of Bill v2.0: the VVUQ-01 and VVUQ-02 engineering record that
%% grounds the findings, drawn from papers/VVUQ-05/update-bill/figures-bill.
%% Common visuals: Figure 6 (gate funnel comparison), Figure 7 (procedure
%% timelines), Figure 8 (humanoid safety figure), Table 5 (external standards map).
%% ============================================================
\clearpage
\phantomsection
\addcontentsline{toc}{section}{Appendix A. Visual Engineering Evidence (Non-Operative)}
\usccrosshead{APPENDIX A. VISUAL ENGINEERING EVIDENCE (NON-OPERATIVE)}

\uscprov{1}{This appendix is non-operative. It visualizes the engineering record
of VVUQ-01 and VVUQ-02 that grounds the findings in section 2 and the gate
discipline codified in section 515D, so a reader can see the evidence the bill
rests on without leaving the document. None of these figures is cited as the
legal basis of an operative clause; only the enacted statutes, in-force rules, and
recognized standards named in the bill carry operative weight
\cite{kawchak2026vvuq01, kawchak2026vvuq02}.}

\figslot{Figure 6.}{The Gate, From One Dimension to Ten}{Side by side: the VVUQ-01
six-candidate funnel that accepted one artifact and blocked five over the verify,
validate, quantify dimensions, and the VVUQ-02 ten-gate funnel that cleared a
surgical humanoid behavior against ten standards-bound gates, three of them
carrying a hard catastrophe predicate.}

\uscprov{1}{The gate is the load-bearing object of the whole record: VVUQ-01 holds
it to a single behavior across three dimensions, and VVUQ-02 scales it to ten
humanoid-specific gates without changing the ACCEPT, BLOCK, ESCALATE semantics.
The verification fraction stays an equality at 1.0 in both, which is why the bill
can require it as an equality \cite{kawchak2026vvuq01, kawchak2026vvuq02}.}

\figslot{Figure 7.}{The Procedure Timelines}{The VVUQ-02 sixty-second eight-phase
Whipple performed by one humanoid with two hands, set over the VVUQ-01 hundred-
and-sixty-eight-day PDAC hybrid pilot, so the single-procedure assurance window
and the full adjuvant pilot window are read on one page.}

\figslot{Figure 8.}{The Humanoid Safety Figure}{A combined view of the
seventy-one-degree-of-freedom kinematic chain, the balance zero-moment-point and
support polygon, and the vessel no-fly proximity bands, the three drawn artifacts
that make the catastrophe gates concrete.}

\uscprov{1}{The safety figure is where embodiment becomes visible: the kinematic
chain shows how many ways a body can move, the support polygon shows the balance
margin a fall would cross, and the no-fly bands show the vascular volume the hard
predicate forbids breaching. These are the dimensions a wrong motion would violate
and a human could not catch in time \cite{isots15066, iec8060127}.}

\figslot{Table 5.}{External Standards Map}{The recognized consensus standards bound
to the gates: ASME V\&V 40, NASA-STD-7009A, IEEE 7009, IEC 80601-2-77, ISO 14971,
ISO/TS 15066, ISO 13482, IEC 60601-1, ISO 13849-1, and ISO 9283, with the gate or
duty each governs.}

\begin{draftbox}{Appendix A - the four visual-evidence artifacts}
\dstep{Figure 6 (gate funnel comparison): render two funnels side by side in the asciifig environment. The VVUQ-01 funnel (6 candidates to 1 accepted, with the per-stage verification, validation, and uncertainty cuts) comes verbatim from papers/VVUQ-05/update-bill/figures-bill/02VVUQ01.md (the "VVUQ Gate Decision Funnel" image spec and the gate decision-surface table). The VVUQ-02 ten-gate funnel (candidates to ACCEPT, with the per-gate V and CV thresholds and the three immediate-catastrophe gates) comes verbatim from papers/VVUQ-05/update-bill/figures-bill/03VVUQ02.md (the "Ten VVUQ Gate Funnel (candidates to ACCEPT)" block). Preserve original spacing on a white background.}
\dstep{Figure 7 (procedure timelines): render the sixty-second eight-phase Whipple timeline from papers/VVUQ-05/update-bill/figures-bill/03VVUQ02.md ("60-Second 8-Phase Whipple Timeline") over the one-hundred-sixty-eight-day PDAC hybrid pilot timeline from papers/VVUQ-05/update-bill/figures-bill/02VVUQ01.md (the PDAC pilot timeline and pdac\_hybrid\_pilot\_timeline.txt block). Keep both as fenced ASCII timelines in the asciifig environment; do not redraw as a chart.}
\dstep{Figure 8 (humanoid safety figure): combine the three drawn artifacts from papers/VVUQ-05/update-bill/figures-bill/03VVUQ02.md - the "H2-Surgical 1.0 Kinematic Chain (71 DOF)", the "Balance ZMP and Support Polygon", and the "Vessel No-Fly Proximity Bands (VVUQ 06)" - into one full-page text figure, each panel labeled, all in the asciifig environment with exact spacing on a white background.}
\dstep{Table 5 (external standards map): build from papers/VVUQ-05/update-bill/figures-bill/03VVUQ02.md (the "Per-gate primary standards", "Cross-cutting standards", and "standards-map" tables) reconciled with the ten-gate schedule in papers/VVUQ-04/final-bill/main.tex (section 515D(c)(4)); columns Standard, What it governs, and Gate or duty; left aligned and ragged right at the body measure (each width prepended with \raggedright\arraybackslash).}
\dstep{Synthesis guidance: these four artifacts are the visual basis for the findings in section 2 and the gate schedule in section 515D, so keep the numbers consistent across them (verification fraction 1.0; the three hard-predicate gates 06, 09, 10; composite mean 93.56; seed 20260525). Use the correlation map in papers/VVUQ-05/update-bill/figures-bill/README.md (Section 4, the cross-file map and the shared evidentiary objects) to confirm each artifact's source and that the 10-gate schedule is the same object cited as evidence in VVUQ-03 and anchored to Title 21 in VVUQ-04.}
\dstep{Formatting: single hyphens only; the section symbol where a codified reference appears; even RaggedRight spacing; keep every figure and table on a white background; prefer full vertical page figures and avoid large empty white bands by sizing each asciifig to its content.}
\end{draftbox}
